\subsection{Setting 1: structured streams, gold dependencies}
\begin{sidewaystable}\centering\footnotesize\setlength{\tabcolsep}{3.5pt}
\caption{Setting 1 (structured streams, gold dependencies, no LLM; 500 streams). Stream-level means with 95\% cluster-bootstrap intervals.}\label{tab:s1}
\begin{tabular}{lcccccccc}
\toprule
System & Warr.\ downstream & Warr.\ direct & Resolvable & Unresolved & Commit rate$\downarrow$ & Unsupp.\ prop.$\downarrow$ & Prov.\ P & Prov.\ R \\
\midrule
\textsc{Palimpsest} & 1.000 {\tiny[1.00,1.00]} & 1.000 {\tiny[1.00,1.00]} & 1.000 {\tiny[1.00,1.00]} & 1.000 {\tiny[1.00,1.00]} & 0.000 {\tiny[0.00,0.00]} & 0.000 {\tiny[0.00,0.00]} & 0.990 {\tiny[0.99,0.99]} & 0.993 {\tiny[0.99,0.99]} \\
\quad no propagation & 0.794 {\tiny[0.78,0.81]} & 1.000 {\tiny[1.00,1.00]} & 0.933 {\tiny[0.93,0.94]} & 0.962 {\tiny[0.96,0.97]} & 0.002 {\tiny[0.00,0.00]} & 0.017 {\tiny[0.01,0.02]} & 0.981 {\tiny[0.98,0.98]} & 0.989 {\tiny[0.99,0.99]} \\
\quad forced commit & 1.000 {\tiny[1.00,1.00]} & 1.000 {\tiny[1.00,1.00]} & 1.000 {\tiny[1.00,1.00]} & 0.042 {\tiny[0.04,0.05]} & 0.776 {\tiny[0.76,0.79]} & 0.090 {\tiny[0.08,0.10]} & 0.992 {\tiny[0.99,0.99]} & 0.775 {\tiny[0.77,0.78]} \\
Log + rule adjudicator & 1.000 {\tiny[1.00,1.00]} & 1.000 {\tiny[1.00,1.00]} & 1.000 {\tiny[1.00,1.00]} & 1.000 {\tiny[1.00,1.00]} & 0.000 {\tiny[0.00,0.00]} & 0.000 {\tiny[0.00,0.00]} & 1.000 {\tiny[1.00,1.00]} & 1.000 {\tiny[1.00,1.00]} \\
Temporal DB & 0.807 {\tiny[0.79,0.82]} & 0.909 {\tiny[0.89,0.92]} & 0.907 {\tiny[0.90,0.91]} & 0.000 {\tiny[0.00,0.00]} & 0.848 {\tiny[0.84,0.86]} & 0.105 {\tiny[0.09,0.12]} & 0.950 {\tiny[0.95,0.95]} & 0.341 {\tiny[0.33,0.35]} \\
Always abstain & 0.000 {\tiny[0.00,0.00]} & 0.000 {\tiny[0.00,0.00]} & 0.000 {\tiny[0.00,0.00]} & 0.393 {\tiny[0.38,0.41]} & 0.000 {\tiny[0.00,0.00]} & 0.075 {\tiny[0.06,0.09]} & 0.507 {\tiny[0.50,0.51]} & 0.739 {\tiny[0.73,0.75]} \\
Always commit & 0.438 {\tiny[0.42,0.46]} & 0.970 {\tiny[0.96,0.98]} & 0.770 {\tiny[0.76,0.78]} & 0.000 {\tiny[0.00,0.00]} & 0.793 {\tiny[0.78,0.81]} & 0.157 {\tiny[0.14,0.17]} & 0.939 {\tiny[0.93,0.94]} & 0.284 {\tiny[0.28,0.29]} \\
\bottomrule
\end{tabular}
\end{sidewaystable}
Table~\ref{tab:s1} reports all systems on 500 streams (\sOneAnswers{} scored answers per system). \pal{} reproduces the oracle on every query, so the incremental, versioned maintenance is correct and the bitemporal slots (as-of valid time, as-of belief time) are answered from stored versions without recomputation. Removing retraction propagation costs $\sOnePropEffect$ points $[\sOnePropLo, \sOnePropHi]$ on warranted downstream queries (paired over \sOnePropN{} streams) and nothing on direct ones---the signature the hypothesis predicts: ordinary belief updating handles the retracted key, only propagation handles its consequences. The forced-commit variant keeps perfect resolvable accuracy while committing on \sOneCommitRate\% of unresolved slots and changing downstream answers after non-warranted events $\sOneCommitUnsup$ points $[\sOneCommitUnsupLo, \sOneCommitUnsupHi]$ more often: premature commitment and unsupported propagation are measured as separate things and move independently of resolvable accuracy. The references split as required (always-abstain: \sOneAbstainRes{} resolvable / \sOneAbstainUnres{} unresolved; always-commit: \sOneCommitRes{} / \sOneCommitUnres{}). The temporal DB, which has valid time but no justifications, loses \sOneTdbWarr{} points on warranted downstream queries and commits on \sOneTdbCommit\% of unresolved slots. Provenance precision and recall of \pal{} are \sOneProvP/\sOneProvR{} against justification paths; the references cite existing but non-supporting ids. These are validation results on a benchmark whose generator and system share an author: they establish that the instruments discriminate, not that H1 holds.

\begin{table}[t]\centering\small
\caption{Setting 1: exact accuracy by query slot (mean over 500 streams).}\label{tab:slots}
\begin{tabular}{lcccccc}
\toprule
Slot & \textsc{Palimpsest} & no propagation & forced commit & Temporal DB & Always commit & Always abstain \\
\midrule
as-of valid time & 1.000 & 1.000 & 0.743 & 0.653 & 0.638 & 0.168 \\
as-of belief time & 1.000 & 1.000 & 0.808 & 0.777 & 0.782 & 0.223 \\
current & 1.000 & 1.000 & 0.723 & 0.675 & 0.701 & 0.242 \\
downstream (derived) & 1.000 & 0.870 & 0.664 & 0.578 & 0.296 & 0.000 \\
reported & 1.000 & 1.000 & 1.000 & 1.000 & 1.000 & 1.000 \\
changed? & 1.000 & 1.000 & 0.932 & 0.919 & 0.919 & 0.068 \\
erroneous? & 1.000 & 1.000 & 0.704 & 0.001 & 0.001 & 0.296 \\
holds at $t$? & 1.000 & 1.000 & 0.877 & 0.771 & 0.734 & 0.224 \\
\bottomrule
\end{tabular}
\end{table}
Table~\ref{tab:slots} breaks accuracy down by slot: the systems without justifications fail on \emph{erroneous?} almost completely (they never hold an error hypothesis), lose most on downstream slots, and answer as-of-belief-time questions from the current state rather than from what was justified then.

\begin{table}[t]\centering\small
\caption{Sensitivity of the Setting 1 gate to the dispute semantics (500 streams). Warranted downstream queries per stream, and the propagation on$-$off difference on them.}\label{tab:sens}
\begin{tabular}{lcccc}
\toprule
Policy & warr./stream & \textsc{Palimpsest} & no propagation & Temporal DB \\
\midrule
P0c (main) & 7.42 & 1.000 {\tiny[1.00,1.00]} & 0.794 {\tiny[0.78,0.81]} & 0.807 {\tiny[0.79,0.82]} \\
literal P0 & 3.38 & 1.000 {\tiny[1.00,1.00]} & 0.638 {\tiny[0.61,0.67]} & 0.918 {\tiny[0.90,0.93]} \\
P1 & 3.49 & 1.000 {\tiny[1.00,1.00]} & 0.688 {\tiny[0.66,0.72]} & 0.866 {\tiny[0.84,0.89]} \\
\bottomrule
\end{tabular}
\end{table}
Table~\ref{tab:sens} shows the sensitivity to the dispute semantics. Under literal P0 only withdrawal resolves, so a stream yields \sOneLitPZeroWarr{} warranted downstream queries and the propagation effect is larger ($+\sOneLitPZeroEffect$ points) because the surviving warranted queries are all retraction-driven; under P1 (class-prioritised dispute with corroboration) \sOnePOneWarr{} per stream and $+\sOnePOneEffect$ points. The qualitative conclusion does not depend on the policy; the size of the usable stratum does.

\subsection{RQ3: calibrated support}
\begin{table}[t]\centering\small
\caption{RQ3, Setting 1 (held-out streams): claim-level calibration of support against the hidden world. Lower is better.}\label{tab:rq3}
\begin{tabular}{lccc}
\toprule
Confidence & Brier & ECE & AURC \\
\midrule
Evidence model, declared priors & 0.182 & 0.118 & 0.138 \\
Evidence model, learned class priors & 0.175 & 0.113 & 0.131 \\
\quad without origin grouping & 0.179 & 0.121 & 0.136 \\
\quad without recency & 0.236 & 0.202 & 0.156 \\
\quad without cue adjustments & 0.183 & 0.128 & 0.136 \\
Class prior only & 0.255 & 0.227 & 0.168 \\
Logistic regression on evidence features & 0.138 & 0.034 & 0.112 \\
Constant 0.5 & 0.250 & 0.496 & 0.268 \\
\bottomrule
\end{tabular}
\end{table}
Support is evaluated at the claim level against the hidden world on 200 held-out streams (\RqN{} candidate values from unresolved and established value slots), with class priors learned on 300 training streams. The evidence model with learned priors is well calibrated (ECE \RqLearnedEce, Brier \RqLearnedBrier) and far better than class priors alone (\RqPriorEce, \RqPriorBrier); recency is the largest single contributor (without it ECE \RqNoRecEce), origin grouping the smallest (\RqNoOrigEce). It is, however, worse than a logistic regression over the same evidence features trained on the same streams (ECE \RqLogregEce, Brier \RqLogregBrier, AURC \RqLogregAurc{} vs.\ \RqLearnedAurc). We report this as a negative result for the hand-designed combination rule: the evidence model's value is its defined meaning and its auditability, not its calibration advantage. H3 proper, against verbalised and sampling-based confidence, is tested in Setting 2 below.

\subsection{Poisoning under an explicit attacker model}
\begin{table}[t]\centering\small
\caption{Poisoning (Setting 1, 60 streams per cell). Attacker injects false assertions at a rate of the honest ones, half with a false change cue; harm = share of value queries answered \emph{established} with the poisoned value; support = mean support the evidence model assigns to poisoned vs.\ honest candidates.}\label{tab:poison}
\begin{tabular}{llcccccc}
\toprule
Channel & Rate & \multicolumn{2}{c}{\textsc{Palimpsest}} & \multicolumn{2}{c}{Temporal DB} & \multicolumn{2}{c}{Support} \\
 & & harm$\downarrow$ & unresolved & harm$\downarrow$ & unresolved & poisoned & honest \\
\midrule
low & 0.01 & 0.003 & 0.27 & 0.003 & 0.03 & 0.21 & 0.67 \\
low & 0.05 & 0.024 & 0.28 & 0.030 & 0.03 & 0.33 & 0.67 \\
low & 0.20 & 0.112 & 0.30 & 0.151 & 0.03 & 0.32 & 0.65 \\
trusted & 0.01 & 0.005 & 0.27 & 0.005 & 0.03 & 0.63 & 0.67 \\
trusted & 0.05 & 0.027 & 0.28 & 0.033 & 0.03 & 0.66 & 0.66 \\
trusted & 0.20 & 0.115 & 0.30 & 0.153 & 0.03 & 0.63 & 0.63 \\
\bottomrule
\end{tabular}
\end{table}
Table~\ref{tab:poison} gives the secondary result. At 20\% injection through a low-trust channel, \pal{} answers \emph{established} with the poisoned value on \PoisHarmPal\% of value queries against \PoisHarmTdb\% for the temporal DB and \PoisHarmCommit\% for always-commit; the difference is smaller than one might expect from a justification-based store because a poisoned report carrying a false change cue is shielded from earlier competitors under A-CUE, so cue resolution is an attack surface. The support layer discriminates poisoned claims through the low channel (mean support \PoisSupLowPois{} vs.\ \PoisSupLowHon{} for honest candidates) but not through a compromised trusted channel (\PoisSupTrPois{} vs.\ \PoisSupTrHon): provenance cannot detect what provenance has endorsed, and only corroboration across origins could. Both effects are what ``measured, not assumed from provenance'' was meant to expose.

\subsection{Setting 2: natural language, automatic extraction}
\begin{sidewaystable}\centering\footnotesize\setlength{\tabcolsep}{3.5pt}
\caption{Setting 2 (natural-language streams, automatic extraction with Claude Haiku 4.5; 100 streams, 99 with a warranted downstream query). Stream-level means with 95\% cluster-bootstrap intervals; cost is measured tokens (input $+8\times$ output) per stream.}\label{tab:s2}
\begin{tabular}{lccccccc}
\toprule
System & Warr.\ downstream & Warr.\ direct & Resolvable & Unresolved & Commit rate$\downarrow$ & Prov.\ P & Cost (tok/stream) \\
\midrule
\textsc{Palimpsest} & 0.956 {\tiny[0.93,0.98]} & 0.996 {\tiny[0.99,1.00]} & 0.976 {\tiny[0.97,0.98]} & 0.834 {\tiny[0.80,0.87]} & 0.080 {\tiny[0.06,0.10]} & 0.978 {\tiny[0.97,0.98]} & 188k \\
\quad no propagation & 0.749 {\tiny[0.70,0.79]} & 0.996 {\tiny[0.99,1.00]} & 0.920 {\tiny[0.91,0.93]} & 0.813 {\tiny[0.78,0.85]} & 0.079 {\tiny[0.06,0.10]} & 0.972 {\tiny[0.97,0.98]} & 188k \\
\quad forced commit & 0.979 {\tiny[0.96,0.99]} & 0.996 {\tiny[0.99,1.00]} & 0.987 {\tiny[0.98,0.99]} & 0.045 {\tiny[0.03,0.06]} & 0.758 {\tiny[0.73,0.79]} & 0.985 {\tiny[0.98,0.99]} & 188k \\
Symbolic replay (rule adjudicator) & 0.956 {\tiny[0.93,0.98]} & 0.996 {\tiny[0.99,1.00]} & 0.976 {\tiny[0.97,0.98]} & 0.834 {\tiny[0.80,0.87]} & 0.080 {\tiny[0.06,0.10]} & 0.989 {\tiny[0.99,0.99]} & 188k \\
Claims log + LLM adjudicator & 0.514 {\tiny[0.46,0.56]} & 0.731 {\tiny[0.69,0.77]} & 0.666 {\tiny[0.64,0.69]} & 0.128 {\tiny[0.11,0.15]} & 0.572 {\tiny[0.53,0.61]} & 0.747 {\tiny[0.73,0.76]} & 532k \\
\quad top-40 entries (original, 30 str.) & 0.419 {\tiny[0.32,0.52]} & 0.731 {\tiny[0.64,0.81]} & 0.614 {\tiny[0.57,0.66]} & 0.103 {\tiny[0.08,0.13]} & 0.504 {\tiny[0.46,0.54]} & 0.758 {\tiny[0.73,0.78]} & 454k \\
Raw-report log + LLM adjudicator (30 str.) & 0.330 {\tiny[0.24,0.43]} & 0.685 {\tiny[0.59,0.78]} & 0.554 {\tiny[0.51,0.60]} & 0.101 {\tiny[0.07,0.13]} & 0.606 {\tiny[0.55,0.66]} & 0.683 {\tiny[0.66,0.71]} & 439k \\
Claims log + LLM adj.\ (amended contract, 30 str.) & 0.400 {\tiny[0.33,0.47]} & 0.703 {\tiny[0.60,0.80]} & 0.607 {\tiny[0.57,0.65]} & 0.134 {\tiny[0.09,0.18]} & 0.533 {\tiny[0.47,0.59]} & 0.714 {\tiny[0.69,0.74]} & 552k \\
\quad top-10 entries & 0.225 {\tiny[0.16,0.29]} & 0.787 {\tiny[0.71,0.86]} & 0.572 {\tiny[0.53,0.62]} & 0.138 {\tiny[0.10,0.17]} & 0.377 {\tiny[0.33,0.43]} & 0.836 {\tiny[0.82,0.85]} & 399k \\
\quad top-20 entries & 0.385 {\tiny[0.30,0.48]} & 0.723 {\tiny[0.63,0.81]} & 0.610 {\tiny[0.57,0.65]} & 0.136 {\tiny[0.10,0.18]} & 0.412 {\tiny[0.36,0.46]} & 0.786 {\tiny[0.77,0.80]} & 424k \\
\quad top-40 entries & 0.429 {\tiny[0.33,0.53]} & 0.739 {\tiny[0.65,0.82]} & 0.628 {\tiny[0.59,0.67]} & 0.133 {\tiny[0.10,0.17]} & 0.471 {\tiny[0.42,0.52]} & 0.748 {\tiny[0.73,0.77]} & 464k \\
Raw-report log + LLM adj.\ (amended, 30 str.) & 0.323 {\tiny[0.23,0.42]} & 0.672 {\tiny[0.59,0.76]} & 0.524 {\tiny[0.48,0.57]} & 0.102 {\tiny[0.07,0.14]} & 0.563 {\tiny[0.51,0.61]} & 0.674 {\tiny[0.65,0.70]} & 448k \\
Temporal DB & 0.779 {\tiny[0.74,0.82]} & 0.885 {\tiny[0.85,0.92]} & 0.895 {\tiny[0.88,0.91]} & 0.000 {\tiny[0.00,0.00]} & 0.837 {\tiny[0.81,0.87]} & 0.961 {\tiny[0.95,0.97]} & 188k \\
Always abstain & 0.000 {\tiny[0.00,0.00]} & 0.000 {\tiny[0.00,0.00]} & 0.000 {\tiny[0.00,0.00]} & 0.380 {\tiny[0.35,0.41]} & 0.000 {\tiny[0.00,0.00]} & 0.542 {\tiny[0.52,0.56]} & 188k \\
Always commit & 0.422 {\tiny[0.38,0.46]} & 0.982 {\tiny[0.97,0.99]} & 0.791 {\tiny[0.77,0.81]} & 0.000 {\tiny[0.00,0.00]} & 0.795 {\tiny[0.77,0.82]} & 0.947 {\tiny[0.94,0.96]} & 188k \\
\bottomrule
\end{tabular}
\end{sidewaystable}
\paragraph{Evaluation set.} 100 streams (13{,}094 reports) were rendered, paraphrased by Claude Haiku 4.5 and passed through the cue-recovery check; after two regenerations 56 lines had lost a stated change and 3 had gained one, and 1{,}343 were flagged by the model-side check as reading a dated statement (``since 3 March'') as a change---flags that cannot alter gold but are counted. Adjudication stripped the lost cues from the frozen labels, changing 274 of 6{,}737 gold entries (4.1\%); the frozen set has \sTwoWarrPerStream{} warranted downstream queries per stream, just under the pre-registered floor of 8 (deviation recorded). The shared extractor (Haiku 4.5) has a slot-error rate of \sTwoExtSlot\%, misses a stated cue on \sTwoExtMissed\% of lines, adds none (\sTwoExtSpur\%), errs on dates on \sTwoExtValid\%, and fails on \sTwoExtFail\%.

\paragraph{H1.} On warranted downstream queries, \pal{} over the extracted claims reaches \sTwoPalWarr{} against \sTwoLogWarr{} for the same claims log adjudicated by the LLM at query time: a paired difference of $+\sTwoHOne$ points $[\sTwoHOneLo, \sTwoHOneHi]$ over \sTwoHOneN{} streams, far beyond the pre-registered 10-point threshold, with the interval excluding zero. The gap is not an artefact of the structured schema: the rule-based adjudicator over the same claims (the upper bound of what query-time reasoning could achieve) reaches \sTwoRuleWarr{}, so the LLM adjudicator loses 44 points to reasoning, not to extraction. Removing propagation costs $\sTwoProp$ points $[\sTwoPropLo, \sTwoPropHi]$, the same effect as in Setting 1. The raw-report log with no extraction at all, the pure query-time alternative, reaches \sTwoNlWarr{} on its 30 streams, and the matched-cost variant that shows the adjudicator only the 40 most recent entries plus the queried key reaches \sTwoTopkWarr{}: less context makes the adjudicator worse, not better.

\paragraph{H4.} The LLM adjudicator commits on \sTwoLogCommit{} of unresolved slots against \sTwoPalCommit{} for \pal, while \pal's resolvable accuracy is \sTwoPalRes{} against \sTwoLogRes: both halves of the joint criterion hold. The adjudicator also changes downstream answers after non-warranted events \sTwoLogUnsup{} of the time (\pal: \sTwoPalUnsup). The forced-commit variant is instructive: on extracted claims it scores \sTwoCommitWarr{} on warranted queries, slightly \emph{above} \pal{} (\sTwoPalWarr), because committing to the minimal-error interpretation often resolves a spurious dispute created by an extraction error in the right direction---but it does so by committing on \sTwoCommitCommit{} of genuinely unresolved slots, which is the trade the contract makes visible.

\paragraph{Cost and the crossover (H2).} Costs are measured tokens. Extraction costs \xE{}k tokens per stream (\$\sTwoPalCost{} at batch prices) and \pal{} pays nothing per query; the claims-log adjudicator pays \xAlog{}k tokens per query on top (\$\sTwoLogCost{} per stream at this query load), the raw-report adjudicator \xAnl{}k per query with no extraction. The crossover ratio at which write-time extraction becomes cheaper than query-time adjudication over raw reports is $r^*=\xRstar$ queries per observation; below it the raw log is cheaper, but at every ratio it is far less accurate, so the accuracy--cost curves do not cross in accuracy. With a shared extractor the claims log is strictly dearer than \pal{} at any query load, and less accurate.

\paragraph{Sign stability with a stronger model.}
\begin{sidewaystable}\centering\footnotesize\setlength{\tabcolsep}{3.5pt}
\caption{Sign-stability subset: 15 Setting 2 streams re-extracted and adjudicated with Claude Sonnet 4.6 on the same frozen evaluation set.}\label{tab:s2strong}
\begin{tabular}{lccccccc}
\toprule
System & Warr.\ downstream & Warr.\ direct & Resolvable & Unresolved & Commit rate$\downarrow$ & Prov.\ P & Cost (tok/stream) \\
\midrule
\textsc{Palimpsest} & 1.000 {\tiny[1.00,1.00]} & 1.000 {\tiny[1.00,1.00]} & 1.000 {\tiny[1.00,1.00]} & 0.816 {\tiny[0.73,0.90]} & 0.090 {\tiny[0.05,0.14]} & 0.987 {\tiny[0.98,1.00]} & 189k \\
\quad no propagation & 0.780 {\tiny[0.64,0.90]} & 1.000 {\tiny[1.00,1.00]} & 0.946 {\tiny[0.91,0.98]} & 0.805 {\tiny[0.72,0.89]} & 0.090 {\tiny[0.05,0.14]} & 0.984 {\tiny[0.97,0.99]} & 189k \\
\quad forced commit & 1.000 {\tiny[1.00,1.00]} & 1.000 {\tiny[1.00,1.00]} & 1.000 {\tiny[1.00,1.00]} & 0.048 {\tiny[0.00,0.12]} & 0.751 {\tiny[0.62,0.87]} & 0.989 {\tiny[0.98,1.00]} & 189k \\
Symbolic replay (rule adjudicator) & 1.000 {\tiny[1.00,1.00]} & 1.000 {\tiny[1.00,1.00]} & 1.000 {\tiny[1.00,1.00]} & 0.816 {\tiny[0.73,0.90]} & 0.090 {\tiny[0.05,0.14]} & 0.997 {\tiny[0.99,1.00]} & 189k \\
Claims log + LLM adjudicator & 0.813 {\tiny[0.67,0.93]} & 0.924 {\tiny[0.85,0.98]} & 0.854 {\tiny[0.80,0.91]} & 0.289 {\tiny[0.19,0.39]} & 0.489 {\tiny[0.36,0.63]} & 0.851 {\tiny[0.82,0.88]} & 718k \\
Temporal DB & 0.777 {\tiny[0.64,0.89]} & 0.844 {\tiny[0.67,0.97]} & 0.908 {\tiny[0.86,0.95]} & 0.000 {\tiny[0.00,0.00]} & 0.833 {\tiny[0.72,0.92]} & 0.957 {\tiny[0.94,0.97]} & 189k \\
Always abstain & 0.000 {\tiny[0.00,0.00]} & 0.000 {\tiny[0.00,0.00]} & 0.000 {\tiny[0.00,0.00]} & 0.399 {\tiny[0.30,0.49]} & 0.000 {\tiny[0.00,0.00]} & 0.531 {\tiny[0.47,0.58]} & 189k \\
Always commit & 0.387 {\tiny[0.29,0.48]} & 0.990 {\tiny[0.97,1.00]} & 0.778 {\tiny[0.72,0.83]} & 0.000 {\tiny[0.00,0.00]} & 0.879 {\tiny[0.84,0.92]} & 0.945 {\tiny[0.92,0.97]} & 189k \\
\bottomrule
\end{tabular}
\end{sidewaystable}
Re-extracting and adjudicating 15 streams with Claude Sonnet 4.6 on the same frozen evaluation set (Table~\ref{tab:s2strong}) lowers the extraction slot-error rate to \sStrongExtSlot\% and lifts the adjudicator to \sStrongLogWarr{} on warranted downstream queries; \pal{} reaches \sStrongPalWarr{} and the paired difference remains positive, $+\sStrongHOne$ points $[\sStrongHOneLo, \sStrongHOneHi]$, with the adjudicator still committing on \sStrongLogCommit{} of unresolved slots. The sign is stable; the size of the gap depends strongly on the adjudicator.

\paragraph{RQ2: extraction noise.}
\begin{table}[t]\centering\small
\caption{RQ2 corruption curve (50 Setting 2 streams, mock extractor with injected error rates). Warranted downstream accuracy; PALIMPSEST resolvable and unresolved accuracy.}\label{tab:noise}
\begin{tabular}{llcccccc}
\toprule
Error type & rate & \textsc{Palimpsest} & no prop. & Temporal DB & Always commit & PAL resolvable & PAL unresolved \\
\midrule
wrong value & 0.00 & 1.000 & 0.801 & 0.799 & 0.423 & 1.000 & 1.000 \\
wrong value & 0.05 & 0.878 & 0.706 & 0.738 & 0.404 & 0.910 & 0.858 \\
wrong value & 0.10 & 0.658 & 0.519 & 0.625 & 0.354 & 0.814 & 0.685 \\
wrong value & 0.20 & 0.449 & 0.332 & 0.388 & 0.257 & 0.652 & 0.515 \\
wrong value & 0.35 & 0.273 & 0.192 & 0.301 & 0.224 & 0.511 & 0.347 \\
change cue dropped & 0.10 & 0.912 & 0.713 & 0.797 & 0.423 & 0.946 & 0.958 \\
change cue dropped & 0.25 & 0.817 & 0.624 & 0.794 & 0.423 & 0.867 & 0.873 \\
change cue dropped & 0.50 & 0.562 & 0.382 & 0.758 & 0.423 & 0.712 & 0.753 \\
change cue added & 0.05 & 1.000 & 0.801 & 0.799 & 0.423 & 1.000 & 0.963 \\
change cue added & 0.15 & 1.000 & 0.815 & 0.799 & 0.423 & 1.000 & 0.871 \\
date dropped & 0.25 & 0.995 & 0.798 & 0.794 & 0.423 & 0.995 & 0.973 \\
date dropped & 0.50 & 0.991 & 0.794 & 0.790 & 0.423 & 0.990 & 0.961 \\
\bottomrule
\end{tabular}
\end{table}
Table~\ref{tab:noise} injects extraction errors at controlled rates with a mock extractor on 50 streams. \pal{} degrades faster than last-write-wins as wrong values accumulate: at 20\% wrong values it scores \nzSlotTwentyPal{} against \nzSlotTwentyTdb{} for the temporal DB, and at 35\% the order reverses (\nzSlotThirtyfivePal{} vs.\ \nzSlotThirtyfiveTdb); when half the stated change cues are dropped, the temporal DB (\nzCueFiftyTdb) is clearly ahead of \pal{} (\nzCueFiftyPal). The mechanism is simple and, we think, the most important limitation of justified-belief maintenance: every extraction error becomes a standing dispute that the semantics keeps open, whereas last-write-wins silently discards it. Dropped dates and added cues cost almost nothing. At the measured error rates of the real extractors (0.3--1.5\%) the effect is small, but it bounds the regime in which the approach pays.

\paragraph{H3: calibration on extracted claims.} For the \RqTwoSupN{} value queries the LLM adjudicator answered as established, the evidence model's support for the asserted value (computed from the extracted evidence) has ECE \RqTwoSupEce{} and Brier \RqTwoSupBrier, worse than the adjudicator's own verbalised confidence (ECE \RqTwoVerbEce, Brier \RqTwoVerbBrier) and than agreement across five samples at temperature 0.7 on a 300-query subset (ECE \RqTwoSampEce, Brier \RqTwoSampBrier); only the risk--coverage areas are close (\RqTwoSupAurc, \RqTwoVerbAurc, \RqTwoSampAurc). H3 is therefore \emph{not} supported in Setting 2: on this population the model's confidence is better calibrated than a provenance-and-recency rule. Together with the Setting 1 comparison against a learned predictor, the evidence model's claim reduces to a defined, auditable meaning of support, not to calibration.

\subsection{Setting 3: streams assembled from Wikidata statement histories}
\begin{sidewaystable}\centering\footnotesize\setlength{\tabcolsep}{3.5pt}
\caption{Setting 3 (transfer: 30 streams assembled from Wikidata statement histories, pseudonymised surface, extraction and adjudication with Claude Haiku 4.5). Descriptive; the warranted-downstream stratum has 24 queries in 13 streams.}\label{tab:s3}
\begin{tabular}{lccccccc}
\toprule
System & Warr.\ downstream & Warr.\ direct & Resolvable & Unresolved & Commit rate$\downarrow$ & Prov.\ P & Cost (tok/stream) \\
\midrule
\textsc{Palimpsest} & 0.923 {\tiny[0.77,1.00]} & 0.976 {\tiny[0.94,1.00]} & 0.978 {\tiny[0.95,1.00]} & 0.909 {\tiny[0.73,1.00]} & 0.000 {\tiny[0.00,0.00]} & 0.892 {\tiny[0.85,0.94]} & 25k \\
\quad no propagation & 0.724 {\tiny[0.49,0.93]} & 0.976 {\tiny[0.94,1.00]} & 0.962 {\tiny[0.93,0.99]} & 0.879 {\tiny[0.70,1.00]} & 0.000 {\tiny[0.00,0.00]} & 0.886 {\tiny[0.84,0.93]} & 25k \\
\quad forced commit & 0.923 {\tiny[0.77,1.00]} & 0.976 {\tiny[0.94,1.00]} & 0.978 {\tiny[0.95,1.00]} & 0.000 {\tiny[0.00,0.00]} & 1.000 {\tiny[1.00,1.00]} & 0.892 {\tiny[0.85,0.94]} & 25k \\
Symbolic replay (rule adjudicator) & 0.923 {\tiny[0.77,1.00]} & 0.976 {\tiny[0.94,1.00]} & 0.978 {\tiny[0.95,1.00]} & 0.909 {\tiny[0.73,1.00]} & 0.000 {\tiny[0.00,0.00]} & 1.000 {\tiny[1.00,1.00]} & 25k \\
Claims log + LLM adjudicator & 0.019 {\tiny[0.00,0.06]} & 0.519 {\tiny[0.38,0.66]} & 0.185 {\tiny[0.14,0.23]} & 0.333 {\tiny[0.09,0.61]} & 0.545 {\tiny[0.27,0.79]} & 0.802 {\tiny[0.75,0.86]} & 88k \\
Raw-report log + LLM adjudicator (30 str.) & 0.077 {\tiny[0.00,0.23]} & 0.451 {\tiny[0.30,0.61]} & 0.181 {\tiny[0.13,0.23]} & 0.182 {\tiny[0.00,0.45]} & 0.576 {\tiny[0.30,0.82]} & 0.731 {\tiny[0.67,0.79]} & 71k \\
Claims log + LLM adj.\ (amended contract, 30 str.) & 0.019 {\tiny[0.00,0.06]} & 0.491 {\tiny[0.34,0.64]} & 0.205 {\tiny[0.15,0.26]} & 0.394 {\tiny[0.15,0.64]} & 0.485 {\tiny[0.21,0.73]} & 0.791 {\tiny[0.73,0.85]} & 92k \\
Raw-report log + LLM adj.\ (amended, 30 str.) & 0.173 {\tiny[0.00,0.38]} & 0.466 {\tiny[0.32,0.61]} & 0.208 {\tiny[0.15,0.26]} & 0.318 {\tiny[0.09,0.59]} & 0.439 {\tiny[0.18,0.70]} & 0.728 {\tiny[0.67,0.79]} & 25k \\
Temporal DB & 0.865 {\tiny[0.65,1.00]} & 0.884 {\tiny[0.78,0.97]} & 0.692 {\tiny[0.64,0.74]} & 0.000 {\tiny[0.00,0.00]} & 1.000 {\tiny[1.00,1.00]} & 0.922 {\tiny[0.88,0.96]} & 25k \\
Always abstain & 0.000 {\tiny[0.00,0.00]} & 0.000 {\tiny[0.00,0.00]} & 0.000 {\tiny[0.00,0.00]} & 0.788 {\tiny[0.55,1.00]} & 0.000 {\tiny[0.00,0.00]} & 0.828 {\tiny[0.78,0.88]} & 25k \\
Always commit & 0.000 {\tiny[0.00,0.00]} & 0.506 {\tiny[0.36,0.65]} & 0.478 {\tiny[0.43,0.53]} & 0.000 {\tiny[0.00,0.00]} & 0.909 {\tiny[0.73,1.00]} & 0.803 {\tiny[0.75,0.86]} & 25k \\
\bottomrule
\end{tabular}
\end{sidewaystable}
Setting 3 is the transfer setting: 30 streams that nobody on the project designed, assembled from the revision histories of 40 Wikidata persons and the 128 organisations recorded as their employers (3,474 sampled revisions; selection, mapping and schema decisions are listed in the released note \texttt{docs/SETTING3.md}). Every edit between sampled revisions becomes an observation: a claim added is an assertion (with a \emph{since} cue when it carries a start qualifier), a value edited in place is a correction, a statement deprecated or removed by the editors is a withdrawal. References give the source class (unreferenced statements are \emph{low}; 331 of 417 assertions). The surface is pseudonymised with the same fictional names as Settings 1--2, so the backbone cannot answer from knowledge of the real person; the timing, conflicts and withdrawals are real. One schema decision matters and is ours: Wikidata's employer property lists every affiliation over a career, mostly undated, so it is declared a stable, multi-valued set (under a single-valued or a changeable reading the derived slot is unresolved almost everywhere: 1 warranted downstream query in 30 streams). The derived slot is therefore the set of headquarters cities of the recorded affiliations, and a warranted downstream query is one where a withdrawal or correction of an affiliation, or a headquarters change, alters that set. The resulting set has 497 observations, 80 withdrawals, \sThreeWarrTotal{} warranted downstream queries in 13 streams (\sThreeWarrPerStream{} per stream), and is reported descriptively, as pre-registered.

Table~\ref{tab:s3} gives the same systems and pipeline as Setting 2 (Haiku 4.5 extraction, slot-error rate \sThreeExtSlot\%). \pal{} over the extracted claims reaches \sThreePalWarr{} on warranted downstream queries, equal to the rule-based adjudicator over the same claims (\sThreeRuleWarr): the ceiling set by extraction. Removing propagation costs $\sThreeProp$ points $[\sThreePropLo, \sThreePropHi]$ over 13 streams, so retraction propagation matters on real histories as it did on generated ones. The LLM adjudicator reaches \sThreeLogWarr{} on the same queries and \sThreeLogRes{} on resolvable ones, a paired difference of $+\sThreeHOne$ points $[\sThreeHOneLo, \sThreeHOneHi]$, and we decompose this gap because most of it is not reasoning in the sense Setting 2 measured. Two output conventions account for a large part: \sThreeEstEmpty{} of the \sThreeEstTotal{} established queries have the empty set as gold (a multi-valued slot whose every reading justifies no member, typically after the last affiliation is withdrawn), which the contract as run did not state and on which the adjudicator always answered \emph{unknown}; and \sThreeHoldsFalse{} yes/no queries are established-false (a value that never held at the asked time, frequent with \emph{since} cues), which the contract does state but the adjudicator again answered \emph{unknown}. Excluding both, its resolvable accuracy is \sThreeLogResNe{} (\pal: \sThreePalResNe). On the warranted downstream queries themselves it answers \emph{unresolved} with the gold set among its alternatives in most cases: it reasons that overlapping affiliations ``may have ended'' although the declared schema marks the attribute stable, i.e.\ it applies its prior about employment to a schema that contradicts it. The temporal DB, which has no such prior, reaches \sThreeTdbWarr{} on warranted downstream queries and commits on every unresolved slot. We draw three conclusions: write-time maintenance transfers to real revision histories (the propagation effect and the equality with the rule adjudicator); the output contract must state the empty-set convention, which we have now added (dated); and the LLM adjudicator's failure here is partly a schema-versus-prior conflict that a pseudonymised, declared schema makes unusually sharp---a limitation of this transfer design as much as of the adjudicator.

\subsection{External check: LongMemEval knowledge updates}
\begin{table}[t]\centering\footnotesize
\caption{LongMemEval knowledge-update items (78; 69 mapped to an extracted attribute). Systems answer one structured slot from claims extracted turn by turn; a Haiku judge compares the answer with the reference. Coverage = share answered as established; accuracy over all items and over answered items. Not comparable to published LongMemEval scores.}\label{tab:lme}
\begin{tabular}{lccc}
\toprule
System & Coverage & Acc.\ (all) & Acc.\ (answered) \\
\midrule
\pal{} (P0c) & 0.32 & 0.167 & 0.400 \\
\pal{} + self-update (exploratory) & 0.74 & 0.474 & 0.621 \\
Temporal DB (last write wins) & 0.88 & 0.526 & 0.594 \\
Always-commit & 0.88 & 0.526 & 0.594 \\
Always-abstain & 0.00 & 0.038 & 0.000 \\
Claims log + LLM adjudicator & 0.56 & 0.359 & 0.614 \\
\bottomrule
\end{tabular}
\end{table}
LongMemEval's knowledge-update items \citep{longmemeval} ask for a user attribute that was stated and later updated across chat sessions. We run the \lmeN{} items of the oracle release through an open-vocabulary adapter shared by all systems: each user turn is extracted into claims about the user (Haiku), attribute keys are canonicalised per item, the question is mapped to one key, and every system answers that slot; a Haiku judge compares the answer with the reference, and on the six items whose reference is ``not enough information'' abstention is credited. These numbers are not comparable to published LongMemEval scores (the systems answer a structured slot, not the question), and the adapter caps all systems: 9 items map to no extracted key.

The result (Table~\ref{tab:lme}) is a scoring-target mismatch rather than a measurement of memory quality, and we report it as such. The benchmark's reference for an update is the user's \emph{latest} statement; nothing in the data verifies it, and the extractor recovered a stated change on 5 of 91 claims on the queried keys. Under our semantics a later uncued value from the same source is a dispute, so \pal{} answers \emph{unresolved} on \lmePalUnres{} of 69 mapped items (coverage \lmePalCov, accuracy \lmePalAcc), where last-write-wins (the temporal DB) answers and is right on 34 of them (accuracy \lmeTdbAcc, coverage \lmeTdbCov); on the items where both commit they agree. The LLM adjudicator abstains on 44\% and reaches \lmeLogAcc. An exploratory variant that reads a later self-report on a single-valued key as a stated change (added after seeing this result, not pre-registered, shown only to make the mechanism visible) recovers \lmePalSuAcc{} at coverage \lmePalSuCov. The benchmark therefore measures the recency convention for single-source self-reports, which is a reasonable convention for that setting and which our semantics deliberately does not assume; the ten items where recency committed and was judged wrong are mostly questions about previous values or about the change itself, which a single current-value slot cannot express, and three presupposition-false questions where recency answers anyway. What LongMemEval does not contain---unreliable sources, hidden world state, withdrawn evidence---is what Settings 1--3 measure. STALE has no public data release, so no further external check was possible.
